% Solutions of chapter 25 -- Text algorithms and compression.
\section*{Chapter \ref{ch:b2:text-algorithms} --- Text Algorithms and Compression}

\begin{solution}{exo:b2:text-algorithms:1}
For the prefixes \texttt{a}, \texttt{aa}, \texttt{aab}, \texttt{aaba}, \texttt{aabaa}, \texttt{aabaaa}, \texttt{aabaaab}
the longest \omterm{def:b2:text-algorithms:border}{borders} are the \omterm{def:b2:finite-automata:word}{empty word}, \texttt{a}, the \omterm{def:b2:finite-automata:word}{empty word}, \texttt{a}, \texttt{aa}, \texttt{aa} and
\texttt{aab}. For $q = 6$, the \omterm{def:b2:text-algorithms:border}{border} \texttt{aa} of \texttt{aabaa} cannot be extended (\texttt{aab} is not a \omterm{def:b2:finite-automata:factor}{suffix} of
\texttt{aabaaa}), but the shorter \omterm{def:b2:text-algorithms:border}{border} \texttt{a} can: \texttt{aa}.

\omoutput{code/bachelor-2/25-text-algorithms/outputs/exo_failure.txt}{The \omterm{def:b2:text-algorithms:border}{failure function} of \texttt{aabaaab}.}
\end{solution}

\begin{solution}{exo:b2:text-algorithms:2}
\omoutput{code/bachelor-2/25-text-algorithms/outputs/exo_kmp_bug.txt}{The shorter version against the expected
\omterm{def:b2:text-algorithms:matching}{occurrences}.}

After $q \gets \pi(q)$ the pattern has moved, and the letter $c$ that caused the mismatch must be compared with
$P[\pi(q)]$; \code{kmp_short} skips it. It misses an \omterm{def:b2:text-algorithms:matching}{occurrence} whenever that letter is the start of the continued match:
on \texttt{aaab}, after \texttt{aa} the third \texttt{a} fails against \texttt{b}, $q$ becomes $1$, and that \texttt{a}
should have made it $2$. When $\pi(q) = 0$ and $c \neq P[0]$, as in the first and last tests, nothing is lost. The
correction is the inner \code{while} loop of \cref{lst:b2:text-algorithms:kmp}, which compares $c$ again until it
matches or $q = 0$.
\end{solution}

\begin{solution}{exo:b2:text-algorithms:3}
If $s$ and $s + d$ are both \omterm{def:b2:text-algorithms:matching}{occurrences} with $0 < d < m$, then $P[d..m-1] = T[s+d..s+m-1] = P[0..m-d-1]$: the pattern
has a \omterm{def:b2:text-algorithms:border}{border} of length $m - d$. So $m - d \leq \pi(m)$, that is $d \geq m - \pi(m)$, and the shifts between are not
\omterm{def:b2:text-algorithms:matching}{occurrences}.

\omcode{code/bachelor-2/25-text-algorithms/listings/exo_bm_all.py}{7}{23}{Boyer--Moore, with the period of the pattern
as the shift after an \omterm{def:b2:text-algorithms:matching}{occurrence}}

\omoutput{code/bachelor-2/25-text-algorithms/outputs/exo_bm_all.txt}{Comparisons before and after the change, with the
same \omterm{def:b2:text-algorithms:matching}{occurrences}.}

The gain is small: \texttt{aba} has period $2$, $a^5$ period $1$ (no gain), and the archive's \omterm{def:b2:text-algorithms:matching}{occurrences} are rare. The
good-suffix rule of Boyer and Moore generalises this shift to every partial match.
\end{solution}

\begin{solution}{exo:b2:text-algorithms:4}
Naive search: on $a^n$ with $a^{m-1}b$, every shift compares $m$ letters, $m(n - m + 1)$ in all; with $ba^{m-1}$, every
shift fails at once, $n - m + 1$. Boyer--Moore compares from the right, so the roles swap. With $a^{m-1}b$, the
\texttt{b} fails against \texttt{a} at $j = m - 1$, with $\lambda(a) = m - 2$: one comparison and a shift of $1$,
$n - m + 1$ in all. With $ba^{m-1}$, the $m - 1$ letters \texttt{a} match and the \texttt{b} fails at $j = 0$, with
$\lambda(a) = m - 1$: the shift is $\max(1, -(m-1)) = 1$, so $m(n - m + 1)$ comparisons, the worst case.
\end{solution}

\begin{solution}{exo:b2:text-algorithms:5}
The encoder emits \texttt{s}, \texttt{h}, \texttt{e}, the space, \texttt{s}, \texttt{e}, \texttt{l}, \texttt{l},
\texttt{s}, then 259 for \texttt{\textvisiblespace s}, \texttt{e}, \texttt{a}, 259 again, 257 for \texttt{he}, 262 for
\texttt{ll} and \texttt{s}; each code adds the word followed by the next byte.

\omoutput{code/bachelor-2/25-text-algorithms/outputs/exo_lzw.txt}{The codes (spaces written \texttt{\_}) and the entries
added.}
\end{solution}

\begin{solution}{exo:b2:text-algorithms:6}
\omoutput{code/bachelor-2/25-text-algorithms/outputs/exo_lzw_bug.txt}{The decoder that looks every code up.}

It fails on \texttt{abababa}: the last code, 258, is the entry the encoder created when it emitted the code before it,
and the decoder, one entry behind, does not have it yet. By the proof of \cref{thm:b2:text-algorithms:lzw}, it fails
exactly on the inputs where the encoder emits a code at the step right after creating it, that is where $w_{k+1} = w_k\,
w_k[0]$: the input contains $x v x v x$ with $x$ a byte, $xv$ the dictionary word $w_k$ read at that point, and $xvx$ a
new entry.
\end{solution}

\begin{solution}{exo:b2:text-algorithms:7}
The walk of \code{words} with a stop as soon as $k$ words are found; the smallest letter is on top of the stack, so the
words come out in alphabetical order.

\omcode{code/bachelor-2/25-text-algorithms/listings/exo_autocomplete.py}{6}{17}{The first $k$ completions, and the
nodes visited}

\omoutput{code/bachelor-2/25-text-algorithms/outputs/exo_autocomplete.txt}{Completions of four prefixes among the 80
book titles; the walk visits a fraction of the subtrie.}
\end{solution}

\begin{solution}{exo:b2:text-algorithms:8}
Here $\pi(q) = q - 1$ for $q \leq m - 1$, since $a^{q-1}$ is the longest \omterm{def:b2:text-algorithms:border}{border} of $a^q$. The first $m - 1$ letters each
match with one comparison: $q = m - 1$. For each letter at $i = m - 1, \dots, n - 2$, the \texttt{b} of the pattern fails
against \texttt{a}, $q$ becomes $m - 2$, the shift $i - (m - 2) \leq n - m$ still fits, and $P[m - 2] = a$ matches: two
comparisons, and $q = m - 1$ again. At $i = n - 1$ the \texttt{b} fails once, and the next shift, $n - m + 1$, no longer
fits. In all $(m - 1) + 2(n - m) + 1 = 2n - m$. The C program confirms it for $m = 100$: $399{,}900$ comparisons for
$n = 200{,}000$.
\end{solution}

\begin{solution}{exo:b2:text-algorithms:9}
\omcode{code/bachelor-2/25-text-algorithms/listings/exo_rk_multi.py}{7}{28}{Rabin--Karp for several patterns: one \omterm{def:b2:text-algorithms:rolling}{rolling
hash}, a dictionary from hash values to patterns}

\omoutput{code/bachelor-2/25-text-algorithms/outputs/exo_rk_multi.txt}{Six words of six letters, each count checked
against \code{str.count}.}

The cost is $\BigO(n)$ for the hashes plus the checks of the windows whose hash is one of the patterns', whatever the
number of patterns, against one pass per pattern for KMP.
\end{solution}

\begin{solution}{exo:b2:text-algorithms:10}
The difference $x - y$ is a non-zero integer with $|x - y| < 256^m = 2^{8m}$. If it had $8m$ distinct prime divisors,
their product, at least $2^{8m}$, would divide it, which is impossible: it has fewer than $8m$. A window $x \neq P$ is a
\omterm{def:b2:text-algorithms:rolling}{spurious hit} for the modulus $q$ when $x \equiv P \pmod q$, that is when $q$ divides $x - P$. With $q$ uniform among the
$K$ smallest primes, this has probability below $8m/K$. Summing over the $n - m + 1$ shifts, the expected number of
\omterm{def:b2:text-algorithms:rolling}{spurious hits} is below $8m(n - m + 1)/K$; with $K$ of the order of $mn$, it is a constant, and the expected cost is
$\BigO(n + m)$ plus the \omterm{def:b2:text-algorithms:matching}{occurrences}' checks. The bound holds for every text, chosen before $q$: the randomness is in
the algorithm, not in the input.
\end{solution}

\begin{solution}{exo:b2:text-algorithms:11}
This is the algorithm of Aho and Corasick. The \omterm{def:b2:text-algorithms:trie}{trie} is built first; the failure links are then set breadth first, so that
the link of a node's parent, which points to a shorter word, is already known, and followed as KMP follows $\pi$. A node
also reports the patterns that end at its failure link.

\omcode{code/bachelor-2/25-text-algorithms/listings/exo_aho.py}{7}{40}{A \omterm{def:b2:text-algorithms:trie}{trie} of patterns with failure links, and a
search that reads the text once}

\omoutput{code/bachelor-2/25-text-algorithms/outputs/exo_aho.txt}{In \texttt{ushers}: \texttt{she}, \texttt{he} and
\texttt{hers}; on the archive, seven words in one pass.}

% ledger: b2:aho-corasick-1975
Alfred Aho and Margaret Corasick published the method in 1975, for searching bibliographies for many keywords at once.
\end{solution}

\begin{solution}{exo:b2:text-algorithms:12}
By induction, the $k$-th code is that of $a^k$ and adds $a^{k+1}$, as long as at least $k + 1$ letters remain: before
it, the dictionary holds $a, a^2, \dots, a^k$ and no longer power of $a$, and the input left starts with $a^{k+1}$, so
the longest dictionary word is $a^k$, followed by an \texttt{a}. After $c$ such codes, $1 + 2 + \dots + c = c(c+1)/2$
letters are coded. The last code takes the $r \leq c + 1$ letters left, which form a dictionary word. The number of codes
is therefore the least $c$ with $c(c+1)/2 \geq n$, and solving $c^2 + c - 2n \geq 0$ gives $c \geq (\sqrt{8n+1} -
1)/2$, so $c = \lceil (\sqrt{8n+1}-1)/2 \rceil$.

\omoutput{code/bachelor-2/25-text-algorithms/outputs/exo_lzw_an.txt}{The number of codes for $n = 1, \dots, 21$, and the
formula checked up to 3000.}
\end{solution}

\begin{solution}{pb:b2:text-algorithms:1}
\textbf{1.} Every line starts with the first eight letters of the pattern, so the naive search compares about nine
letters at each of the 1,503 line starts, instead of one; a pattern whose first letter is rare fails at once almost
everywhere.

\textbf{2.} At most shifts the first letter differs, one comparison; the shifts where a prefix matches add a few more.
The cost is $n$ plus the total length of the partial matches, small on ordinary text.

\textbf{3.} About $87{,}299 / 6{,}865 \approx 13$ letters per comparison, and nearly every shift makes one comparison,
so the average jump is about 13. The longest jumps, $21$, come from letters absent from the pattern: digits, most
capitals, punctuation.

\textbf{4.} A date such as \texttt{2025-12-12}: its last letters are digits and the bad character is almost always a
digit, which occurs near the end of the pattern, so the jumps are of one to three letters.

\textbf{5.} It reads each window from its right end and jumps ahead, so it needs the $m$ letters of a window before
deciding anything, and it reads text letters out of order. KMP reads each letter once, in order, and knows after each
one whether an \omterm{def:b2:text-algorithms:matching}{occurrence} ends there.

\textbf{6.} KMP, or the automaton: its whole state is $q$, saved at the end of the day. The other two must keep the last
$m - 1$ letters of yesterday.

\textbf{7.} $\pi$ is $0$ everywhere except $\pi(5) = 1$ and $\pi(14) = 1$: the only \omterm{def:b2:text-algorithms:border}{borders} are \texttt{r}, for
\texttt{retur} and \texttt{returned Under}.

\textbf{8.} After a mismatch KMP falls back to $q = 0$ or $1$ and compares the same letter again, which is what the naive
search does at the next shift when the match was short; both make about one comparison per letter plus the partial
matches.

\textbf{9.} $2n - m = 2 \times 87{,}299 - 18 = 174{,}580$. The observed $91{,}581$ is about $n$: the bound counts a
second comparison for every letter, which only a pattern with long \omterm{def:b2:text-algorithms:border}{borders} on a matching text forces.

\textbf{10.} Naive: $100 \times (10^6 - 99) = 99{,}990{,}100$; KMP: $2 \times 10^6 - 100 = 1{,}999{,}900$, as printed.

\textbf{11.} For a fixed pattern, $m(n - m + 1)$ is linear in $n$ with slope $m$. It is not linear in the size of the
input $n + m$: for $m = n/2$ it is about $n^2/4$.

\textbf{12.} States $0$ to $3$: $\delta(0, a) = 1$, $\delta(1, a) = 2$, $\delta(2, a) = 2$, $\delta(2, b) = 3$,
$\delta(3, a) = 1$, and every other transition leads to $0$.

\textbf{13.} Each node below the root stands for one distinct prefix of a name; a letter is saved each time a name
reaches a prefix that an earlier name had already created, such as the \texttt{W}, \texttt{e}, \texttt{n},
\texttt{\textvisiblespace} shared by the two names starting with \texttt{Wen}. The $152 = 751 - 599$ saved letters are
the letters of shared prefixes, counted once for each extra name sharing them.

\textbf{14.} Six nodes, the root and one per letter, whatever the number of names; each step is one lookup in a small
table of children.

\textbf{15.} A sorted list finds the first name with the prefix by halving, $\BigO(|u| \log N)$ letter comparisons for
$N$ names, then reads the $r$ answers in order. The \omterm{def:b2:text-algorithms:trie}{trie} needs $|u|$ steps and the walk of the subtrie. Both are fast for
60 names; the \omterm{def:b2:text-algorithms:trie}{trie} does not depend on $N$.

\textbf{16.} Each node other than the root has two children or more or ends a name, and there are 60 names, so at most
$59$ branching nodes besides the 60 ends: at most $2 \times 60 = 120$ nodes with the root.

\textbf{17.} Keep the node reached by the prefix typed so far, and move down one edge when a letter is added, or back up
when one is deleted.

\textbf{18.} A \omterm{def:b2:text-algorithms:trie}{trie} of the names written as ``last first'' as well doubles the letters inserted, at most $2 \times 751 +
1$ nodes, and lets a prefix of either name be completed; two separate \omterm{def:b2:text-algorithms:trie}{tries} are the same size.

\textbf{19.} The codes of \texttt{a}, \texttt{aa}, \texttt{aaa}, \texttt{aaaa}, that is $97, 256, 257, 258$; the entries
$256 = \texttt{aa}$, $257 = \texttt{aaa}$, $258 = \texttt{aaaa}$, and nothing after the last code, which ends the
input.

\textbf{20.} As in \cref{exo:b2:text-algorithms:12}. For $n = 55 = 10 \cdot 11/2$, $c = 10$; for $5050 = 100 \cdot
101/2$, $c = 100$; one more letter needs one more code, $101$.

\textbf{21.} The dictionary learns the archive's phrases as it reads: at the start it holds pairs of bytes; after
thousands of lines it holds \texttt{borrowed\textvisiblespace}, names, titles and the frame of the weather lines, and
each code covers more bytes.

\textbf{22.} $11{,}383 \times 2 = 22{,}766$ bytes, against $18{,}095$ with codes of 9 to 14 bits; the largest code,
$256 + 11{,}382$, needs 14 bits.

\textbf{23.} A Huffman code of bytes depends only on the bytes' frequencies, which are about the same in any long part
of the archive: the ratio is $8$ divided by the average code length, about 4.9 bits here, and it cannot use the
repeated words.

\textbf{24.} LZW on $a^n$: exactly $\lceil (\sqrt{8n+1}-1)/2 \rceil$ codes; on the archive, $11{,}383$ codes for
$87{,}299$ bytes, $7.7$ bytes per code and a ratio of $4.82$ with growing code widths, against $1.63$ for Huffman and
$8.63$ for zlib.
\end{solution}
